%% Question bank chapter file - generated by bank_tools.py build (do not edit while the bank is still being built:
%% it is regenerated from the per-test data in Question Bank\data). Every question is one bankq block:
%%   \begin{bankq}[id=..., source={...}, paper=.., qno=.., type=..., subtopic={...}, answer={...}, answer_src=key|solved,
%%                 working={...}, difficulty=1-10, diff_parts={K S I C F}, diff_note={...},
%%                 group=<passage id or empty>, also_in={duplicate copies in other tests}]
%% difficulty = min(10, K+S+I+C+F-1): K concepts 1-3, S steps 1-4, I insight 0-2, C computation 0-2, F format 0-1
%%   ...question LaTeX, options, figures (\bfig / \bfigi -> img/)...
%%   \end{bankq}
%% Passages shared by several questions are bankpassage blocks; their questions carry group=<passage id>.
\documentclass[10pt]{article}
\usepackage[a4paper,margin=1.8cm]{geometry}
\usepackage[T1]{fontenc}
\usepackage{amsmath,amssymb,mathtools}
\usepackage{graphicx}
\usepackage[export]{adjustbox}
\usepackage{xcolor,needspace,keyval}
\usepackage{enumitem}
\usepackage{booktabs,array,tabularx}
\usepackage[hidelinks]{hyperref}
\setlength{\parindent}{0pt}
\setlength{\parskip}{3pt}
\makeatletter
\define@key{bq}{id}{\def\bq@id{#1}}\define@key{bq}{source}{\def\bq@src{#1}}\define@key{bq}{paper}{\def\bq@paper{#1}}
\define@key{bq}{qno}{\def\bq@qno{#1}}\define@key{bq}{type}{\def\bq@type{#1}}\define@key{bq}{subtopic}{\def\bq@sub{#1}}
\define@key{bq}{answer}{\def\bq@ans{#1}}\define@key{bq}{answer_src}{\def\bq@asrc{#1}}\define@key{bq}{working}{\def\bq@work{#1}}
\define@key{bq}{group}{\def\bq@group{#1}}\define@key{bq}{also_in}{\def\bq@also{#1}}
\define@key{bq}{difficulty}{\def\bq@diff{#1}}\define@key{bq}{diff_parts}{\def\bq@dparts{#1}}\define@key{bq}{diff_note}{\def\bq@dnote{#1}}
\newenvironment{bankq}[1][]{\def\bq@work{}\def\bq@also{}\def\bq@group{}\def\bq@diff{?}\def\bq@dparts{}\def\bq@dnote{}\setkeys{bq}{#1}%
  \par\needspace{4\baselineskip}\medskip\noindent{\small\textbf{\bq@id}\hfill\color{gray}[\bq@src, Paper \bq@paper, Q\bq@qno]
  \bq@type\ $\cdot$ \bq@sub\ $\cdot$ \textbf{D\bq@diff}}\par\nopagebreak\relax}
  {\par{\small\color{teal}Answer: \textbf{\bq@ans}\ifx\bq@work\@empty\else\ (solved: \bq@work)\fi
  \ifx\bq@also\@empty\else\par\color{gray}Also in: \bq@also\fi}\par}
\define@key{bp}{id}{\def\bp@id{#1}}\define@key{bp}{source}{\def\bp@src{#1}}\define@key{bp}{questions}{\def\bp@qs{#1}}
\newenvironment{bankpassage}[1][]{\setkeys{bp}{#1}\par\needspace{4\baselineskip}\bigskip\noindent
  {\small\textbf{Passage \bp@id}\hfill\color{gray}[\bp@src; questions \bp@qs]}\par\nopagebreak\relax\itshape}{\par}
\makeatother
\newcommand{\bfig}[2]{\par\smallskip{\centering\includegraphics[width=#1pt,max width=0.62\linewidth,max height=0.32\textheight]{img/#2}\par}\smallskip}
\newcommand{\bfigi}[2]{\includegraphics[width=#1pt,max width=\linewidth,max height=0.32\textheight,valign=c]{img/#2}}
\newcommand{\bmissing}[1]{\textit{[figure missing in the source conversion: \detokenize{#1}]}}
\newenvironment{qoptions}{\begin{itemize}[leftmargin=1.8em,itemsep=1pt,topsep=2pt,label={}]}{\end{itemize}}
\newcommand{\opt}[2]{\item[(#1)] #2}
\providecommand{\fwstart}{}\providecommand{\fwend}{}\providecommand{\qfullsplit}{}
\newcolumntype{L}[1]{>{\raggedright\arraybackslash}p{#1}}
\newcolumntype{C}[1]{>{\centering\arraybackslash}p{#1}}
\newcolumntype{R}[1]{>{\raggedleft\arraybackslash}p{#1}}
\providecommand{\sectionhdr}[1]{}\providecommand{\subsectionhdr}[1]{}


\begin{document}
{\Large\bfseries Biomolecules --- Question Bank}\par\smallskip
{\small Chemistry chapter C25. Questions from all converted test papers, grouped by source series. Answers: printed key where it matches the paper, otherwise solved (answer\_src=solved, with working). Exact repeats across papers appear once, with also\_in listing the other copies.}\par\medskip\hrule

\section*{Score-I papers}
\begin{bankq}[id=SI12-P2-CHEM-S2-Q01, source={Score-I - Test 12 (09 Mar 2022)}, paper=2, qno=1, type=numerical, subtopic={Amino acids}, answer={4}, answer_src=solved, working={Essential: phenylalanine, arginine, histidine, tryptophan}, difficulty=2, diff_parts={K1 S1 I0 C0 F1}, diff_note={concepts: essential amino acids | steps: recall the essential ones among the eight}, also_in={TMAS - Test 04 (19 Nov 2023) P2 Q3}]
Identify total number of Essential amino acids\\
(a) Glycine \quad (b) Alanine \quad (c) Phenyl alanine \quad (d) Arginine\\
(e) Cysteine \quad (f) Histidine \quad (g) Proline \quad (h) Tryptophan
\end{bankq}

\section*{Score Advanced papers}
\begin{bankq}[id=SA01-P1-CHEM-S2-Q09, source={Score Advanced - Test 01 (22 Sep 2020)}, paper=1, qno=9, type=multi, subtopic={Carbohydrates}, answer={A,B,C}, answer_src=key, difficulty=7, diff_parts={K2 S2 I2 C1 F1}, diff_note={concepts: osazone formation, carbohydrates | steps: reaction of reducing sugars; carbons involved; compare osazones}]
Consider the following reaction :\par
(D) $-$ glucose $\xrightarrow[\text{cat.H}^+]{\text{Ph--NH--NH}_2\ (3\ \text{equiv})}$ X\par
Among the following, the compound(s) whose osazone derivative(s) will have the same melting point as that of X is/are
\begin{qoptions}
\opt{A}{\bfigi{71.5}{SA01_p1_chem_s2_q09_A.png}}
\opt{B}{\bfigi{66.7}{SA01_p1_chem_s2_q09_B.png}}
\opt{C}{\bfigi{66.5}{SA01_p1_chem_s2_q09_C.png}}
\opt{D}{\bfigi{66.5}{SA01_p1_chem_s2_q09_D.png}}
\end{qoptions}
\end{bankq}
\begin{bankq}[id=SA01-P1-CHEM-S2-Q12, source={Score Advanced - Test 01 (22 Sep 2020)}, paper=1, qno=12, type=multi, subtopic={Amino acids}, answer={B,C,D}, answer_src=key, difficulty=3, diff_parts={K1 S1 I1 C0 F1}, diff_note={concepts: electrophoresis of amino acids | steps: charge at the given pH; direction of movement}]
You have a mixture of three amino acids : A(pI = 3), B(pI = 5) and C(pI = 10), where pI represents their isoelectric points. Under electrophoresis at pH = 8, the movement of each component of the mixture is considered. Which of the following options are incorrect regarding movement acids under electrophoresis :
\begin{qoptions}
\opt{A}{A and B to anode; C to cathode}
\opt{B}{A to anode; B stationary; C to cathode}
\opt{C}{A to cathode ; B stationary; C to anode}
\opt{D}{A to anode; B and C to cathode}
\end{qoptions}
\end{bankq}
\begin{bankq}[id=SA01-P1-CHEM-S3-Q06, source={Score Advanced - Test 01 (22 Sep 2020)}, paper=1, qno=6, type=numerical, subtopic={Carbohydrates}, answer={1.75}, answer_src=key, difficulty=8, diff_parts={K2 S2 I2 C2 F1}, diff_note={concepts: mutarotation of sugars, carbohydrates | steps: specific rotation; equilibrium mixture; composition}]
Specific rotations of freshly prepared aqueous solutions of \textbf{I} and \textbf{II} are +112 and +18, respectively. On standing the optical rotation of aqueous solution of \textbf{I} slowly decreases give a final value of +52.7 due to equilibration with \textbf{II}. Under this state of equilibrium, what the ratio \textbf{II : I}?\par
\bfig{398.0}{SA01_p1_chem_s3_q06_1.png}
\end{bankq}
\begin{bankq}[id=SA01-P2-CHEM-S3-Q08, source={Score Advanced - Test 01 (22 Sep 2020)}, paper=2, qno=8, type=numerical, subtopic={Amino acids}, answer={9.74}, answer_src=key, difficulty=2, diff_parts={K1 S1 I0 C0 F1}, diff_note={concepts: isoelectric point of lysine | steps: relevant pKa values; average}, also_in={Score Advanced - Test 09 (29 Apr 2023) P2 Q5}]
Calculate the isoelectric point ($p$I) of lysine. Give the pK$_\mathrm{a}$ of $\alpha$-NH$_3$ is 8.95, pK$_\mathrm{a}$ side chain NH$_3$ is 10.53 and pK$_\mathrm{a}$ of $\alpha$-COOH is 2.18.
\end{bankq}
\begin{bankq}[id=SA03-P1-CHEM-S1-Q07, source={Score Advanced - Test 03 (27 Jun 2021)}, paper=1, qno=7, type=multi, subtopic={Amino acids}, answer={B,C,D}, answer_src=solved, working={Phenylalanine, valine and threonine are essential}, difficulty=2, diff_parts={K1 S1 I0 C0 F1}, diff_note={concepts: essential amino acids | steps: recall}]
Which Amino acid(s) is/are essential amino acid -
\begin{qoptions}
\opt{A}{Glycine}
\opt{B}{Phenyl alanine}
\opt{C}{Valine}
\opt{D}{Threonine}
\end{qoptions}
\end{bankq}
\begin{bankq}[id=SA03-P1-CHEM-S2-Q04, source={Score Advanced - Test 03 (27 Jun 2021)}, paper=1, qno=4, type=single, subtopic={Carbohydrates}, answer={A}, answer_src=solved, working={beta-glucose; polysaccharide; non reducing}, difficulty=3, diff_parts={K2 S2 I0 C0 F0}, diff_note={concepts: properties of cellulose, carbohydrates | steps: monomer unit; type; nature}]
Which of the following combination is correct for Cellulose ?
\begin{qoptions}
\opt{A}{ii , 2 , 4}
\opt{B}{i , 2 , 4}
\opt{C}{iv , 4}
\opt{D}{ii , 2 , 3}
\end{qoptions}
\end{bankq}
\begin{bankq}[id=SA03-P1-CHEM-S2-Q05, source={Score Advanced - Test 03 (27 Jun 2021)}, paper=1, qno=5, type=single, subtopic={Carbohydrates}, answer={D}, answer_src=solved, working={Lactose is glucose and beta-galactose, not glucopyranose and fructofuranose}, difficulty=2, diff_parts={K2 S1 I0 C0 F0}, diff_note={concepts: wrongly matched monomer, carbohydrates | steps: monomer of each carbohydrate; compare}]
Which of the following is not correctly match with their monomer ?
\begin{qoptions}
\opt{A}{P , ii}
\opt{B}{Q , iii}
\opt{C}{R , i}
\opt{D}{S , iii}
\end{qoptions}
\end{bankq}
\begin{bankq}[id=SA03-P1-CHEM-S2-Q06, source={Score Advanced - Test 03 (27 Jun 2021)}, paper=1, qno=6, type=single, subtopic={Carbohydrates}, answer={D}, answer_src=solved, working={Starch}, difficulty=1, diff_parts={K1 S1 I0 C0 F0}, diff_note={concepts: non reducing sugar | steps: which sugar has no free anomeric group}]
Which of the following is non reducing sugar ?
\begin{qoptions}
\opt{A}{$\alpha$-Glucose}
\opt{B}{$\beta$-Glucose}
\opt{C}{Galactose}
\opt{D}{Starch}
\end{qoptions}
\end{bankq}
\begin{bankq}[id=SA03-P2-CHEM-S1-Q12, source={Score Advanced - Test 03 (27 Jun 2021)}, paper=2, qno=12, type=single, subtopic={Carbohydrates}, answer={B}, answer_src=solved, working={sucrose}, difficulty=1, diff_parts={K1 S1 I0 C0 F0}, diff_note={concepts: non reducing sugar | steps: sucrose}]
Which of the following is non reducing sugar ?
\begin{qoptions}
\opt{A}{Maltose}
\opt{B}{Sucrose}
\opt{C}{Fructose}
\opt{D}{Lactose}
\end{qoptions}
\end{bankq}
\begin{bankq}[id=SA03-P2-CHEM-S2-Q03, source={Score Advanced - Test 03 (27 Jun 2021)}, paper=2, qno=3, type=single, subtopic={Amino acids}, answer={11}, answer_src=solved, working={C5H11NO2}, difficulty=1, diff_parts={K1 S1 I0 C0 F0}, diff_note={concepts: hydrogen atoms in an essential amino acid | steps: identify valine; formula}]
Which of the following is essential amino acid? Give total number of hydrogen atoms present in it.\par
(a) Proline \quad (b) Glutamine \quad (c) Valine \quad (d) Serine
\end{bankq}
\begin{bankq}[id=SA05-P1-CHEM-S1-Q14, source={Score Advanced - Test 05 (15 Jul 2021)}, paper=1, qno=14, type=single, subtopic={Carbohydrates}, answer={B}, answer_src=key, difficulty=6, diff_parts={K3 S1 I2 C1 F0}, diff_note={concepts: hexose giving two hexahydric alcohols, carbohydrates, epimers reduce to the same alcohol | steps: epimers reduce to the same alcohol; the hexose must be a ketose of a glucose epimer}]
X is a hexose. On catalytic reduction X gives two hexahydric alcohol Y and Z. Y can be obtained from D(+) glucose by reducing with NaBH$_4$. X can be :
\begin{qoptions}
\opt{A}{\bfigi{52.7}{SA05_p1_chem_s1_q14_A.png}}
\opt{B}{\bfigi{58.5}{SA05_p1_chem_s1_q14_B.png}}
\opt{C}{\bfigi{58.5}{SA05_p1_chem_s1_q14_C.png}}
\opt{D}{\bfigi{56.5}{SA05_p1_chem_s1_q14_D.png}}
\end{qoptions}
\end{bankq}
\begin{bankq}[id=SA07-P2-CHEM-S3-Q06, source={Score Advanced - Test 07 (16 Apr 2023)}, paper=2, qno=6, type=numerical, subtopic={Carbohydrates}, answer={5}, answer_src=key, difficulty=8, diff_parts={K3 S1 I2 C2 F1}, diff_note={concepts: reducing sugars and acetylation, carbohydrates, identify reducing sugars | steps: identify reducing sugars; count acetyl groups}]
\bfigi{54.1}{SA07_p2_chem_s3_q06_A.png}\hfill\bfigi{54.1}{SA07_p2_chem_s3_q06_B.png}\par
\bfigi{61.5}{SA07_p2_chem_s3_q06_C.png}\par
\bfigi{115.3}{SA07_p2_chem_s3_q06_D.png}\par
\bfigi{201.3}{SA07_p2_chem_s3_q06_5.png}\par
\bfigi{200.8}{SA07_p2_chem_s3_q06_6.png}\par
Calculate the number of reducing sugar = x\par
Calculate the number of sugar in which molecular weight increases by 210, whenever it undergo acetylation reaction with acetic anhydride = y\par
So x -- y =
\end{bankq}
\begin{bankq}[id=SA08-P1-CHEM-S1-Q03, source={Score Advanced - Test 08 (22 Apr 2023)}, paper=1, qno=3, type=single, subtopic={Proteins}, answer={C}, answer_src=key, difficulty=3, diff_parts={K1 S1 I1 C1 F0}, diff_note={concepts: enzyme properties | steps: enzymes work at physiological pH}]
Which of the following information is not correct with respect to enzymes and their catalytic action.
\begin{qoptions}
\opt{A}{Enzymes are complex nitrogenous compounds which form colloidal solution in water.}
\opt{B}{Enzymes are protein molecules of high molecular mass where as co-enzymes are non protein used for enhancing catalytic action of enzyme.}
\opt{C}{The rate of enzyme catalysed reaction is higher at high pH of around 12.}
\opt{D}{The enzyme urease cannot catalyze hydrolysis of any amide other than urea.}
\end{qoptions}
\end{bankq}
\begin{bankq}[id=SA13-P1-CHEM-S2-Q03, source={Score Advanced - Test 13 (13 May 2023)}, paper=1, qno=3, type=single, subtopic={Proteins}, answer={4}, answer_src=key, difficulty=7, diff_parts={K3 S1 I2 C2 F0}, diff_note={concepts: peptide linkages and amino acid residues in a polypeptide, proteins, number of residues gives the linkages | steps: number of residues gives the linkages}]
1 mole of a polypeptide of hydrolysis produced 4 moles of glycine, 2 moles of alanine, 1 mole of phenyl alanine and 1 mole of leucine.\par
If x = number of peptide linkages in polypeptide molecule\par
y = possible number of polypeptide if all identical amino acids remain together in all polypeptides.\par
z = Mass of water (gm) gained by 1 mole of polypeptide on complete hydrolysis.\par
Calculate : $\dfrac{yz}{108x}$\par
\end{bankq}
\begin{bankq}[id=SA14-P1-CHEM-S3-Q05, source={Score Advanced - Test 14 (16 May 2023)}, paper=1, qno=5, type=numerical, subtopic={Proteins}, answer={18}, answer_src=key, difficulty=8, diff_parts={K3 S1 I2 C2 F1}, diff_note={concepts: sequence count from a pentapeptide, proteins, hydrolysis data | steps: hydrolysis data; arrangement of residues}]
A pentapeptide has $-$COOH group on tyrosine. This produces (Gly), valine (Val), Phenylalanine (Phe), Alanine (Ala) and Tyrosine (Tyr) on complete hydrolysis. For this pentapeptide, the number of possible sequences (primary structures) with $-$NH$_2$ group attached to a chiral centre is \_\_\_\_\_\_\_\_\_\par
\end{bankq}
\begin{bankq}[id=SA17-P2-CHEM-S3-Q07, source={Score Advanced - Test 17 (25 May 2023)}, paper=2, qno=7, type=numerical, subtopic={Amino acids}, answer={0.375}, answer_src=key, difficulty=7, diff_parts={K3 S1 I2 C1 F1}, diff_note={concepts: essential amino acids and phosphodiester carbons, count essential amino acids, 3 and 5 positions of the sugar | steps: count essential amino acids; 3 and 5 positions of the sugar}]
(a) Leucine, Lysine, Serine, Cysteine, Proline, Threonine, Arginine, Histidine\par
Number of Essential amino acids in above mentioned are X.\par
(b) In phosphodiester linkage of a Dinucleotide, carbon (Y) of one sugar and carbon number (Z) of other sugar are joined by phosphate. Then X $\div$ (Y + Z) is :
\end{bankq}
\begin{bankq}[id=SA18-P1-CHEM-S1-Q05, source={Score Advanced - Test 18 (27 May 2023)}, paper=1, qno=5, type=multi, subtopic={Carbohydrates}, answer={B,C}, answer_src=key, difficulty=8, diff_parts={K3 S1 I2 C2 F1}, diff_note={concepts: structure of a polymeric fructan, carbohydrates, identify glycosidic linkage and sugar units | steps: identify glycosidic linkage and sugar units}]
The polymeric structure shown in the diagram contains:\par
\bfig{192.8}{SA18_p1_chem_s1_q05_1.png}
\begin{qoptions}
\opt{A}{At least one $\alpha$-D-glucopyranose unit}
\opt{B}{Poly-$\beta$-D-fructofuranoside moity}
\opt{C}{(2 $\rightarrow$ 1) $\beta$-glycosidic linkage}
\opt{D}{At least one $\beta$-D-fructofuranose unit}
\end{qoptions}
\end{bankq}
\begin{bankq}[id=SA19-P1-CHEM-S3-Q07, source={Score Advanced - Test 19 (04 Apr 2024)}, paper=1, qno=7, type=numerical, subtopic={Carbohydrates}, answer={4}, answer_src=key, difficulty=8, diff_parts={K3 S1 I2 C2 F1}, diff_note={concepts: stereoisomers of an aldose that give a meso alditol, carbohydrates, symmetry of the reduced product | steps: symmetry of the reduced product}]
Observe the following reaction and find out that how many number of reactant stereoisomers can be reduced to optically inactive meso product\par
\bfig{148.2}{SA19_p1_chem_s3_q07_1.png}
\end{bankq}
\begin{bankq}[id=SA20-P1-CHEM-S3-Q05, source={Score Advanced - Test 20 (06 Apr 2024)}, paper=1, qno=5, type=numerical, subtopic={Proteins}, answer={18}, answer_src=key, difficulty=9, diff_parts={K3 S2 I2 C2 F1}, diff_note={concepts: number of pentapeptide sequences with a given terminal, proteins, fix the C-terminal tyrosine | steps: fix the C-terminal tyrosine; N-terminal not glycine; arrange the rest}]
A pentapeptide has $-$COOH group on tyrosine. This produces (Gly), valine (Val), Phenylalanine (Phe), Alanine (Ala) and Tyrosine (Tyr) on complete hydrolysis. For this pentapeptide, the number of possible sequences (primary structures) with $-$NH$_2$ group attached to a chiral centre is \underline{\hspace{1.5cm}}
\end{bankq}
\begin{bankq}[id=SA22-P1-CHEM-S1-Q05, source={Score Advanced - Test 22 (10 Apr 2024)}, paper=1, qno=5, type=single, subtopic={Amino acids}, answer={C}, answer_src=key, difficulty=4, diff_parts={K2 S1 I1 C1 F0}, diff_note={concepts: compound that does not exist as a zwitterion, amino acids | steps: amino acids versus amino compounds that are not acids}]
Which of the following does not exist as a Zwitter ion?
\begin{qoptions}
\opt{A}{Glycine}
\opt{B}{Para Sulphanilic acid}
\opt{C}{\bfigi{125.8}{SA22_p1_chem_s1_q05_C.png}}
\opt{D}{Alanine}
\end{qoptions}
\end{bankq}
\begin{bankq}[id=SA23-P1-CHEM-S2-Q07, source={Score Advanced - Test 23 (14 Apr 2024)}, paper=1, qno=7, type=multi, subtopic={Carbohydrates}, answer={A,B,D}, answer_src=key, difficulty=8, diff_parts={K3 S1 I2 C2 F1}, diff_note={concepts: identification of two aldohexoses from osazone and oxidation, carbohydrates, same osazone means C2 epimers | steps: same osazone means C2 epimers; nitric acid oxidation to a meso aldaric acid}]
On treatment with phenylhydrazine, aldohexoses \textbf{P} and \textbf{Q} give the same osazone. On treatment with warm nitric acid, \textbf{P} gives an optically inactive aldaric acid, but sugar \textbf{Q} gives an optically active aldaric acid. Sugars \textbf{P} and \textbf{Q} are both degraded to aldopentose \textbf{R}, which gives an optically active aldaric acid on treatment with nitric acid. Aldopentose \textbf{R} is degraded to aldotetrose \textbf{S}, which gives optically active tartaric acid when it is treated with nitric acid. Aldotetrose \textbf{S} is degraded to D(+)-glyceraldehyde. Identify correct structures of sugars ?
\begin{qoptions}
\opt{A}{\bfigi{59.3}{SA23_p1_chem_s2_q07_A.png}}
\opt{B}{\bfigi{59.6}{SA23_p1_chem_s2_q07_B.png}}
\opt{C}{\bfigi{57.0}{SA23_p1_chem_s2_q07_C.png}}
\opt{D}{\bfigi{58.2}{SA23_p1_chem_s2_q07_D.png}}
\end{qoptions}
\end{bankq}
\begin{bankq}[id=SA25-P2-CHEM-S4-Q10, source={Score Advanced - Test 25 (21 Apr 2024)}, paper=2, qno=10, type=single, subtopic={Amino acids}, answer={2}, answer_src=key, difficulty=7, diff_parts={K3 S1 I2 C2 F0}, diff_note={concepts: positively charged amino acids at pH 7, basic amino acids lysine arginine histidine | steps: basic amino acids lysine arginine histidine}, also_in={Score Advanced - Test 26 (21 Apr 2024) P2 Q10}]
How many amino acids are positively charged at pH = 7.\par
(i) Asparagine \hfill (ii) Valine \hfill (iii) Aspartic acid\par
(iv) Glycine \hfill (v) Lysine \hfill (vi) Glutamine\par
(vii) Glutamic acid \hfill (viii) Arginine \hfill (ix) Methionine
\end{bankq}
\begin{bankq}[id=SA28-P2-CHEM-S3-Q01, source={Score Advanced - Test 28 (27 Apr 2024)}, paper=2, qno=1, type=matching, subtopic={Carbohydrates}, answer={P$\to$2; Q$\to$1; R$\to$4; S$\to$3}, answer_src=solved, working={Sucralose is a derivative of the disaccharide sucrose (2), aspartame is a dipeptide derivative (1), galactaric acid is the oxidation product of the aldohexose galactose (4), and xylitol is the reduction product of an aldopentose (3).}, difficulty=9, diff_parts={K3 S2 I2 C2 F1}, diff_note={concepts: classification of sucralose aspartame galactaric acid and xylitol, carbohydrates, sucralose from sucrose | steps: sucralose from sucrose; aspartame from dipeptide; aldaric acid from aldohexose; alditol from aldopentose}]
\par
{\footnotesize\setlength{\tabcolsep}{1.5pt}\noindent\begin{tabular}{|c|L{4.2cm}|c|L{2.6cm}|}
\hline
 & \multicolumn{1}{>{\centering\arraybackslash}p{4.2cm}|}{List-I (Compound)} & & \multicolumn{1}{>{\centering\arraybackslash}p{2.6cm}|}{List-II (Characteristic)} \\ \hline
(P) & \bfigi{95.1}{SA28_p2_chem_s3_pre_1.png} & (1) & A derivative of dipeptide \\ \hline
(Q) & \bfigi{98.9}{SA28_p2_chem_s3_pre_2.png} & (2) & A derivative of disaccharide \\ \hline
(R) & \bfigi{60.8}{SA28_p2_chem_s3_pre_3.png} & (3) & Reduction product of an aldopentose \\ \hline
(S) & \bfigi{40.0}{SA28_p2_chem_s3_pre_4.png} & (4) & Oxidation product of aldohexose \\ \hline
\end{tabular}}\par
Match each process given in List-I with one or more effect(s) in List-II.\par
\end{bankq}
\begin{bankq}[id=SA30-P1-CHEM-S3-Q04, source={Score Advanced - Test 30 (04 May 2024)}, paper=1, qno=4, type=numerical, subtopic={Amino acids}, answer={9}, answer_src=key, difficulty=8, diff_parts={K3 S1 I2 C2 F1}, diff_note={concepts: standard problem in this chapter, amino acids, apply the relevant results | steps: apply the relevant results; compute}, also_in={Score Advanced - Test 31 (04 May 2024) P1 Q4}]
If pKa values for protonated Lysine molecule are pKa$_1$ = 2.18 ; pKa$_2$ = 8.95 ; pKa$_3$ = 10.79 respectively then , Calculate vlue of (x $-$ 1) where x is iso-electric point of lysine.\par
\bfig{212.8}{SA30_p1_chem_s3_q04_1.png}
Answer should be nearest integral value
\end{bankq}
\begin{bankpassage}[id=SA32-P2-CHEM-S3-G5-6, source={Score Advanced - Test 32 (08 May 2024)}, questions={5-6}]
The isoelectric point (P$_\mathrm I$) is the P$_\mathrm H$ at which the amino acid exists only as a dipolar ion with net charge zero. Structure of lysine and aspartic acid are $\mathrm{C\overset{\ominus}{O}O\,CH\left(\overset{+}{N}H_3\right)(CH_2)_4NH_2}$ and $\mathrm{CO\overset{\ominus}{O}\,CH\left(\overset{+}{N}H_3\right)CH_2COOH}$ respectively. The P$_{\mathrm{Ka}_1}$, P$_{\mathrm{Ka}_2}$ and P$_{\mathrm{Ka}_3}$ of the dictation lysine are 2.18, 8.95 and 10.53 respectively. The P$_{\mathrm{Ka}_1}$, P$_{\mathrm{Ka}_2}$ and P$_{\mathrm{Ka}_3}$ of the cation of aspartic acid are 1.88, 3.65 and 9.60 respectively.
\end{bankpassage}
\begin{bankq}[id=SA32-P2-CHEM-S3-Q05, source={Score Advanced - Test 32 (08 May 2024)}, paper=2, qno=5, type=numerical, subtopic={Amino acids}, answer={9.74}, answer_src=key, difficulty=8, diff_parts={K3 S1 I2 C2 F1}, diff_note={concepts: standard problem in this chapter, amino acids, apply the relevant results | steps: apply the relevant results; compute}, group=SA32-P2-CHEM-S3-G5-6]
Isoelectric point lysine is\_\_\_\_\_\_\_\_\_\_.
\end{bankq}
\begin{bankq}[id=SA32-P2-CHEM-S3-Q06, source={Score Advanced - Test 32 (08 May 2024)}, paper=2, qno=6, type=numerical, subtopic={Amino acids}, answer={2.77}, answer_src=key, difficulty=8, diff_parts={K3 S1 I2 C2 F1}, diff_note={concepts: standard problem in this chapter, amino acids, apply the relevant results | steps: apply the relevant results; compute}, group=SA32-P2-CHEM-S3-G5-6]
Isoelectric point of Aspartic acid is \_\_\_\_\_\_\_\_\_\_.
\end{bankq}
\begin{bankq}[id=SA33-P1-CHEM-S2-Q06, source={Score Advanced - Test 33 (12 May 2024)}, paper=1, qno=6, type=multi, subtopic={Carbohydrates}, answer={A,B,C}, answer_src=key, difficulty=8, diff_parts={K3 S1 I2 C2 F1}, diff_note={concepts: standard problem in this chapter, carbohydrates, apply the relevant results | steps: apply the relevant results; compute}, also_in={Score Advanced - Test 37 (19 May 2024) P1 Q6}]
The correct statement/s about the pyranose form of a sugar (X) given below is/are:\par
\bfig{73.0}{SA33_p1_chem_s2_q06_1.png}
\begin{qoptions}
\opt{A}{It exists in two anomeric pyranose forms}
\opt{B}{It reacts with Tollen's reagent to give a silver mirror}
\opt{C}{The penta-O-methyl derivative of (X) is non reducing}
\opt{D}{It resists reduction with aqueous sodium borohydride}
\end{qoptions}
\end{bankq}
\begin{bankq}[id=SA35-P2-CHEM-S3-Q07, source={Score Advanced - Test 35 (15 May 2024)}, paper=2, qno=7, type=numerical, subtopic={Amino acids}, answer={0.375}, answer_src=key, difficulty=7, diff_parts={K3 S1 I2 C1 F1}, diff_note={concepts: essential amino acids and phosphodiester carbons, count essential amino acids, 3 and 5 positions of the sugar | steps: count essential amino acids; 3 and 5 positions of the sugar}]
(a)\ \ Leucine , Lysine , Serine, Cysteine , Proline, Threonine , Arginine , Histidine\par
\hspace*{1.2em}Number of Essential amino acids in above mentioned are X.\par
(b)\ \ In phosphodiester linkage of a Dinucleotide, carbon number (Y) of one sugar and carbon number (Z) of other sugar are joined by phosphate. Then X $\div$ (Y + Z)\ \ is :
\end{bankq}
\begin{bankpassage}[id=SA36-P2-CHEM-S3-G1-2, source={Score Advanced - Test 36 (19 May 2024)}, questions={1-2}]
Consider the following process of folding/unfolding of a tripeptide segment in a large protein. The free energy change ($\Delta$G) will depend on the interaction of the unfolded tripeptide with the solvent (water) and with the rest of the protein in the folded state (see below). Assume that the tripeptide is made up of one non-polar (hydrophobic, shaded) and two polar (hydrophilic; unshaded) residues. Use the following $\Delta$G data of formation of interaction between residue with water or protein :\par
i) a non-polar residue and the solvent (water) $\Delta$G = +8 kJ mol$^{-1}$\par
ii) a non-polar residue and the rest of the protein $\Delta$G = $-$4 kJ mol$^{-1}$\par
iii) a polar residue and the solvent (water) $\Delta$G = $-$16 kJ mol$^{-1}$\par
iv) a polar residue and the rest of the protein $\Delta$G = $-$14 kJ mol$^{-1}$\par
\bfig{201.3}{SA36_p2_chem_s3_para1_2_1.png}
\end{bankpassage}
\begin{bankq}[id=SA36-P2-CHEM-S3-Q01, source={Score Advanced - Test 36 (19 May 2024)}, paper=2, qno=1, type=numerical, subtopic={Proteins}, answer={8}, answer_src=key, difficulty=8, diff_parts={K3 S1 I2 C2 F1}, diff_note={concepts: standard problem in this chapter, proteins, apply the relevant results | steps: apply the relevant results; compute}, group=SA36-P2-CHEM-S3-G1-2]
What is the magnitude of $\Delta$G for the folding of the tripeptide segment in KJ/mol?
\end{bankq}
\begin{bankq}[id=SA36-P2-CHEM-S3-Q02, source={Score Advanced - Test 36 (19 May 2024)}, paper=2, qno=2, type=numerical, subtopic={Proteins}, answer={6}, answer_src=key, difficulty=8, diff_parts={K3 S1 I2 C2 F1}, diff_note={concepts: standard problem in this chapter, proteins, apply the relevant results | steps: apply the relevant results; compute}, group=SA36-P2-CHEM-S3-G1-2]
What is the magnitude of $\Delta$G for the folding of the tripeptide segment if all the three residues are polar in KJ/mol?
\end{bankq}
\begin{bankpassage}[id=SA37-P2-CHEM-S3-G1-2, source={Score Advanced - Test 37 (19 May 2024)}, questions={1-2}]
Consider the following process of folding/unfolding of a tripeptide segment in a large protein. The free energy change ($\Delta$G) will depend on the interaction of the unfolded tripeptide with the solvent (water) and with the rest of the protein in the folded state (see below). Assume that the tripeptide is made up of one non-polar (hydrophobic, shaded) and two polar (hydrophilic; unshaded) residues. Use the following $\Delta$G data of formation of interaction between residue with water or protein :\par
i) a non-polar residue and the solvent (water) $\Delta$G = +8 kJ mol$^{-1}$\par
ii) a non-polar residue and the rest of the protein $\Delta$G = $-$4 kJ mol$^{-1}$\par
iii) a polar residue and the solvent (water) $\Delta$G = $-$16 kJ mol$^{-1}$\par
iv) a polar residue and the rest of the protein $\Delta$G = $-$14 kJ mol$^{-1}$\par
\bfig{201.3}{SA37_p2_chem_s3_para1_2_1.png}
\end{bankpassage}
\begin{bankq}[id=SA37-P2-CHEM-S3-Q01, source={Score Advanced - Test 37 (19 May 2024)}, paper=2, qno=1, type=numerical, subtopic={Proteins}, answer={8}, answer_src=key, difficulty=8, diff_parts={K3 S1 I2 C2 F1}, diff_note={concepts: standard problem in this chapter, proteins, apply the relevant results | steps: apply the relevant results; compute}, group=SA37-P2-CHEM-S3-G1-2]
What is the magnitude of $\Delta$G for the folding of the tripeptide segment in KJ/mol?
\end{bankq}
\begin{bankq}[id=SA37-P2-CHEM-S3-Q02, source={Score Advanced - Test 37 (19 May 2024)}, paper=2, qno=2, type=numerical, subtopic={Proteins}, answer={6}, answer_src=key, difficulty=8, diff_parts={K3 S1 I2 C2 F1}, diff_note={concepts: standard problem in this chapter, proteins, apply the relevant results | steps: apply the relevant results; compute}, group=SA37-P2-CHEM-S3-G1-2]
What is the magnitude of $\Delta$G for the folding of the tripeptide segment if all the three residues are polar in KJ/mol?
\end{bankq}

\section*{Enthusiast Pre-Score papers}
\begin{bankq}[id=EPS01-P1-CHEM-S3-Q06, source={Enthusiast Pre-Score - Test 01 (04 Nov 2020)}, paper=1, qno=6, type=numerical, subtopic={Proteins}, answer={4.00}, answer_src=key, difficulty=4, diff_parts={K1 S2 I1 C0 F1}, diff_note={concepts: biuret test | steps: peptide linkages; select compounds; count}]
Identify compounds give Biurate test?\par
(i) Cellulose \quad (ii) Protien \quad (iii) Nylon-6 \quad (iv) Starch\par
(v)\par\bfig{90.2}{EPS01_p1_chem_s3_q06_1.png}\par(vi) H$_2$N$-$gly$-$Ala$-$Cys$-$COOH
\end{bankq}
\begin{bankq}[id=EPS01-P2-CHEM-S1-Q01, source={Enthusiast Pre-Score - Test 01 (04 Nov 2020)}, paper=2, qno=1, type=multi, subtopic={Carbohydrates}, answer={B,D}, answer_src=key, difficulty=3, diff_parts={K1 S2 I0 C0 F1}, diff_note={concepts: statements about glucose | steps: Schiff test; epimers; oxidation; mutarotation}]
Following statement(s) is/are true for D-glucose :
\begin{qoptions}
\opt{A}{It shows Schiff's reagent test}
\opt{B}{It is C-2 epimer of D-mannose}
\opt{C}{Oxidise with Br$_2$ / H$_2$O to form saccharic acid}
\opt{D}{Shows mutarotation.}
\end{qoptions}
\end{bankq}
\begin{bankq}[id=EPS01-P2-CHEM-S1-Q02, source={Enthusiast Pre-Score - Test 01 (04 Nov 2020)}, paper=2, qno=2, type=multi, subtopic={Amino acids}, answer={B,C,D}, answer_src=key, difficulty=2, diff_parts={K1 S1 I0 C0 F1}, diff_note={concepts: chiral amino acids | steps: which amino acids are chiral}]
Optically active amino acid is :
\begin{qoptions}
\opt{A}{Glycine}
\opt{B}{Alanine}
\opt{C}{Valine}
\opt{D}{Phenyl alanine}
\end{qoptions}
\end{bankq}
\begin{bankq}[id=EPS01-P2-CHEM-S3-Q03, source={Enthusiast Pre-Score - Test 01 (04 Nov 2020)}, paper=2, qno=3, type=single, subtopic={Carbohydrates}, answer={3}, answer_src=key, difficulty=3, diff_parts={K1 S2 I1 C0 F0}, diff_note={concepts: reducing sugars | steps: free anomeric groups; identify; count}]
The number of reducing sugars among the following is \underline{\hspace{0.8cm}}\par
\bfig{387.1}{EPS01_p2_chem_s3_q03_1.png}
\end{bankq}

\section*{TAAS papers}
\begin{bankq}[id=TAAS02-P2-CHEM-S4-Q07, source={TAAS - Test 02 (30 Jul 2023)}, paper=2, qno=7, type=numerical, subtopic={Carbohydrates}, answer={10}, answer_src=key, difficulty=5, diff_parts={K1 S2 I1 C1 F1}, diff_note={concepts: functional groups in carbohydrate derivatives | steps: structures P, Q, R, S; count hydroxyl and ketone groups; add}]
\bfig{136.4}{TAAS02_p2_chem_s4_q07_1.png}
Number of $-$OH groups in (P) $=$ ($x$)

Number of $-$OH group in (Q) $=$ ($y$)

Number of alcohol group in (R) $=$ ($z$)

Number of ketone groups in (S) $=$ ($k$)

the value of $x+y+z+k$ is?
\end{bankq}

\section*{TMAS papers}
\begin{bankq}[id=TMAS04-P1-CHEM-S3-Q04, source={TMAS - Test 04 (19 Nov 2023)}, paper=1, qno=4, type=numerical, subtopic={Amino acids}, answer={5}, answer_src=solved, working={Tyr-Asp-Lys: at pH 2 charge $+2$ (both amines protonated); at pH 6 it is $0$ ($+1+1-1-1$); at pH 11 about $-3$ (phenoxide, side-chain carboxylate, C-terminus). Sum $=2+0+3=5$ (pKa values approximate)}, difficulty=6, diff_parts={K1 S2 I2 C1 F1}, diff_note={concepts: ionisation of a peptide | steps: ionisable groups; charge at each pH; sum}]
The structure of a peptide is given below.
\bfig{154.3}{TMAS04_p1_chem_s3_q04_1.png}
If the absolute values of the net charge of the peptide at pH $=2$, pH $=6$, and pH $=11$ are $|z_1|$, $|z_2|$ and $|z_3|$, respectively, then what is $|z_1|+|z_2|+|z_3|$ ?
\end{bankq}
\begin{bankq}[id=TMAS04-P2-CHEM-S2-Q03, source={TMAS - Test 04 (19 Nov 2023)}, paper=2, qno=3, type=numerical, subtopic={Proteins}, answer={6.00}, answer_src=key, difficulty=3, diff_parts={K1 S1 I0 C1 F1}, diff_note={concepts: sequences of a tripeptide | steps: number of arrangements of three}]
A tripeptide is composed equally of L-tyrosine, L-glycine and L-valine (one molecule of each). How many number of tripeptides can be obtained?
\end{bankq}
\begin{bankq}[id=TMAS04-P2-CHEM-S2-Q04, source={TMAS - Test 04 (19 Nov 2023)}, paper=2, qno=4, type=numerical, subtopic={Amino acids}, answer={10.00}, answer_src=key, difficulty=3, diff_parts={K1 S1 I0 C1 F1}, diff_note={concepts: essential amino acids | steps: recall the count}]
How many essential amino acids are present in the living system?
\end{bankq}

\section*{TNAS papers}
\begin{bankq}[id=TNAS05-P1-CHEM-S1-Q19, source={TNAS - Test 05 (09 Jul 2023)}, paper=1, qno=19, type=single, subtopic={Carbohydrates}, answer={B}, answer_src=key, difficulty=7, diff_parts={K2 S3 I1 C2 F0}, diff_note={concepts: reduction of a sugar with red phosphorus and HI, glucose structure | steps: product Q; hydroxyl groups; Tollens test; stereocentres; monochlorination count}]
L$-$glucose $\xrightarrow[\Delta]{\mathrm{Red\ P\ +\ HI}}$ Q\par
Choose the correct option?
\begin{qoptions}
\opt{A}{`Q' contains five hydroxyl groups}
\opt{B}{L-glucose can give positive Tollen's test.}
\opt{C}{Structure of glucose have total 32 stereo isomers.}
\opt{D}{Monochlorination of ``Q'' gives ``20'' monochloro products.}
\end{qoptions}
\end{bankq}

\section*{Star Review (Other) papers}
\begin{bankq}[id=STAR03-P1-CHEM-S1-Q02, source={Star Review (Other) - Test 03 (04 Dec 2023)}, paper=1, qno=2, type=single, subtopic={Carbohydrates}, answer={D}, answer_src=key, difficulty=1, diff_parts={K1 S1 I0 C0 F0}, diff_note={concepts: reducing disaccharides | steps: identify the disaccharides; which ones are non reducing}]
Which of the following is not reducing disaccharide ?
\begin{qoptions}
\opt{A}{Starch}
\opt{B}{Lactose}
\opt{C}{Maltose}
\opt{D}{Sucrose}
\end{qoptions}
\end{bankq}
\begin{bankpassage}[id=STAR03-P2-CHEM-S3-G1-2, source={Star Review (Other) - Test 03 (04 Dec 2023)}, questions={1-2}]
Consider the following reactions\par
$\mathrm{D-(+)-glucose}\xrightarrow{\mathrm{HNO_3}}[\mathrm P]$\par
$\mathrm{D-(+)-glucose}\xrightarrow{\mathrm{Br_2/H_2O}}[\mathrm Q]$\par
The number of stereo isomers possible for $P$ and $Q$ (including $P$ \& $Q$), respectively, are $x$ and $y$.
\end{bankpassage}
\begin{bankq}[id=STAR03-P2-CHEM-S3-Q01, source={Star Review (Other) - Test 03 (04 Dec 2023)}, paper=2, qno=1, type=numerical, subtopic={Carbohydrates}, answer={10.00}, answer_src=key, difficulty=7, diff_parts={K1 S3 I2 C1 F1}, diff_note={concepts: stereoisomers of the glucose oxidation product | steps: oxidation of both ends; structure; stereocentres; symmetry; count}, group=STAR03-P2-CHEM-S3-G1-2]
The value of $x$ is \underline{\hspace{0.8cm}}.
\end{bankq}
\begin{bankq}[id=STAR03-P2-CHEM-S3-Q02, source={Star Review (Other) - Test 03 (04 Dec 2023)}, paper=2, qno=2, type=numerical, subtopic={Carbohydrates}, answer={16.00}, answer_src=key, difficulty=6, diff_parts={K1 S2 I2 C1 F1}, diff_note={concepts: stereoisomers of gluconic acid | steps: oxidation of the aldehyde; stereocentres; count}, group=STAR03-P2-CHEM-S3-G1-2]
The value of $y$ is \underline{\hspace{0.8cm}}.
\end{bankq}

\section*{Nurture (General) papers}
\begin{bankq}[id=NUR01-P1-CHEM-S1-Q01, source={Nurture (General) - Test 01 (12 Jun 2021)}, paper=1, qno=1, type=multi, subtopic={Carbohydrates}, answer={B,C}, answer_src=solved, working={Glucose and fructose are functional isomers; sucrose has one glycosidic link; fructose reacts with Tollens (tautomerism)}, difficulty=3, diff_parts={K1 S2 I0 C0 F1}, diff_note={concepts: statements about sugars | steps: isomers; Tollens test; ring size; glycosidic links}]
Which of the following statement(s) is/are correct ?
\begin{qoptions}
\opt{A}{Glucose \& Fructose are optical isomers.}
\opt{B}{Fructose can be oxidised by tollens reagent.}
\opt{C}{Monosaccharides containing six membered ring can be called pyranose sugar.}
\opt{D}{Sucrose is a disaccharide and it contains two glycosidic linkages.}
\end{qoptions}
\end{bankq}
\end{document}
